\section{Experimental design}\label{sec:experimental-design}

\subsection{Primary benchmark: OpenML-CC18}
\label{sec:exp-datasets}
The primary doctoral benchmark is the OpenML-CC18 curated
classification suite of \citet{bischl2021openmlbenchmark}
(\texttt{suite\_id}~$=$~99), containing 72 standardized classification
tasks served by the OpenML platform \citep{vanschoren2014openml}. The
suite was constructed by filtering OpenML for tasks with verified
licences, well-defined train/test splits, dataset sizes between
500 and 100\,000 rows (with a small number of larger exceptions),
$\geq\!2$ classes, and reasonable class balance, and is widely used
as a fair common ground for HPO and AutoML evaluation. We use it as
\emph{a task-based benchmark, not a dataset-based one}: a few CC18
tasks share the same underlying OpenML dataset under different target
columns or train/test splits, so identity is by
\texttt{openml\_task\_id} and never by dataset name alone. The
materialized task / dataset registry, suite metadata, and a
binary / multiclass / categorical / imbalance / size-bucket coverage
report are committed under
\texttt{benchmarks/doctoral/openml\_cc18/} and are produced
deterministically by \texttt{scripts/import\_openml\_cc18.py}.

\paragraph{Coverage of the panel.} The 72 tasks split as 35 binary
and 37 multiclass; 24 tasks have at least one categorical feature;
19 tasks have a class-imbalance ratio $\geq\!5$. By size, 18 tasks
have $\leq\!1\,000$ rows, 43 have $1\,001$--$30\,000$, and 11 have
$>\!30\,000$. The full per-task table (rows, features, $n$ classes,
class distribution, license, OpenML URL) lives in
\texttt{benchmarks/doctoral/openml\_cc18/tasks.csv}; we do not
reproduce it here because the suite is canonical and externally
specified.

\paragraph{Internal smoke panel (not part of the benchmark).} The
12-dataset panel that drove earlier integration testing
(\texttt{magic}, \texttt{breast\_cancer}, \texttt{pima\_diabetes},
\texttt{spambase}, \texttt{adult}, \texttt{bank\_marketing},
\texttt{credit\_card\_default}, \texttt{german\_credit},
\texttt{wine\_quality}, \texttt{dry\_bean}, \texttt{mushroom},
\texttt{phishing}) is preserved as a development / smoke / runtime
profiling fixture under
\texttt{benchmarks/doctoral/internal\_smoke\_panel/} and is
\emph{not} reported as part of the doctoral benchmark.

\subsection{Algorithms}
\label{sec:exp-algorithms}
We compare three GBDT families: XGBoost \citep{chen2016xgboost},
LightGBM \citep{ke2017lightgbm}, and CatBoost
\citep{prokhorenkova2018catboost}. Categorical features are passed
natively to CatBoost; for XGBoost and LightGBM the headline run uses
label-encoded ints for cross-method parity with HPO baselines that
do not understand categorical columns natively (this default is
toggleable per algorithm and is logged in the per-replica manifest).
The proposed method and every baseline tune the same per-algorithm
hyperparameter list under the same bounds, listed in
Table~\ref{tab:hpo-spaces}.
\todoblock{Confirm the LightGBM / CatBoost search spaces. Initial
plan: keep the seven-factor XGBoost space for parity with the
dissertation, and use analogues for LightGBM (\texttt{num\_leaves},
\texttt{min\_data\_in\_leaf}) and CatBoost (\texttt{depth},
\texttt{l2\_leaf\_reg}).}

\begin{table}[t]
  \centering
  \caption{Per-algorithm hyperparameter search spaces (tentative).}
  \label{tab:hpo-spaces}
  \begin{tabular}{lll}
    \toprule
    Algorithm & Hyperparameter & Range \\
    \midrule
    \multirow{7}{*}{XGBoost}
      & \texttt{subsample}         & $[0.05, 1.00]$ \\
      & \texttt{colsample\_bytree} & $[0.05, 1.00]$ \\
      & \texttt{colsample\_bylevel}& $[0.05, 1.00]$ \\
      & \texttt{learning\_rate}    & $[0.01, 0.30]$ \\
      & \texttt{max\_depth}        & $\{3,\dots,18\}$ \\
      & \texttt{gamma}             & $[0.05, 5.00]$ \\
      & \texttt{n\_estimators}     & $\{50,\dots,700\}$ \\
    \midrule
    \multirow{7}{*}{LightGBM} & \multicolumn{2}{l}{\todoblock{fill in.}}\\
    \midrule
    \multirow{7}{*}{CatBoost} & \multicolumn{2}{l}{\todoblock{fill in.}}\\
    \bottomrule
  \end{tabular}
\end{table}

\subsection{Comparative protocol and method matrix}
\label{sec:exp-protocol}
The methods entering the doctoral benchmark are frozen by the
machine-readable matrix at
\texttt{benchmarks/doctoral/openml\_cc18/method\_matrix.csv}; the
authoritative narrative specification is
\texttt{docs/COMPARATIVE\_PROTOCOL.md}. The matrix is the single
source of truth for shard generation: the script that materializes
the SQLite job matrix reads its method list from that CSV and never
hardcodes method names. Table~\ref{tab:methods} summarizes the
matrix. Methods carry a \emph{primary / subset / ablation /
literature\_only} status and an \emph{objective mode} flag (single-
versus multi-objective).

\begin{table}[t]
  \centering
  \caption{Methods participating in the OpenML-CC18 doctoral
  benchmark. ``Status'' is one of \emph{primary} (runs on all 72
  tasks at every stage), \emph{subset} (runs only on the CC18 subset
  defined in \texttt{benchmarks/doctoral/openml\_cc18/comparative\_protocol.md}),
  \emph{ablation} (joins from stage 2 onward), or \emph{literature}
  (cited as context, not benchmarked). Multi-fidelity rows share a
  budget-equivalence rule based on total boosting iterations rather
  than configuration count. The full machine-readable specification
  is \texttt{benchmarks/doctoral/openml\_cc18/method\_matrix.csv}.}
  \label{tab:methods}
  \small
  \begin{tabular}{lllll}
    \toprule
    Method & Family & Status & Objective & Reference \\
    \midrule
    \texttt{default\_gbdt}                & control          & primary    & single & --- \\
    \texttt{random\_search}               & classical        & primary    & single & \citep{bergstra2012random} \\
    \texttt{tpe\_optuna}                  & classical        & primary    & single & \citep{bergstra2011tpe,akiba2019optuna} \\
    \texttt{smac3}                        & SMAC             & primary    & single & \citep{hutter2011smac,lindauer2022smac3} \\
    \texttt{hyperband\_or\_asha}          & multi-fidelity   & primary    & single & \citep{li2017hyperband,li2020asha} \\
    \texttt{bohb}                         & multi-fidelity   & primary    & single & \citep{falkner2018bohb} \\
    \texttt{dehb}                         & multi-fidelity   & primary    & single & \citep{awad2021dehb} \\
    \texttt{nsga2}                        & evolutionary MOO & primary    & multi  & \citep{deb2002nsga2,blank2020pymoo} \\
    \texttt{motpe}                        & MOO              & primary    & multi  & \citep{ozaki2020motpe} \\
    \texttt{parego}                       & MOO              & subset     & multi  & \citep{knowles2006parego} \\
    \texttt{doe\_rsm\_vrf\_true\_nbi}     & proposed         & primary    & multi  & this work \\
    \texttt{\ldots\_no\_mbpa}             & ablation         & ablation   & multi  & this work \\
    \texttt{legacy\_weighted\_sum}        & ablation         & ablation   & multi  & \citep{ribeiro2026dissertation} \\
    \texttt{flaml\_optional}              & AutoML context   & literature & single & \citep{wang2021flaml} \\
    \bottomrule
  \end{tabular}
\end{table}

\paragraph{Replicas and CV.} Each (task, algorithm, method) cell is
repeated for $R{=}30$ replicas, staged $1\!\to\!5\!\to\!10\!\to\!30$
through the four stages \texttt{stage0\_replica\_001},
\texttt{stage1\_topup\_to\_005}, \texttt{stage2\_topup\_to\_010},
\texttt{stage3\_topup\_to\_030}. Each replica fits the OpenML
task-defined CV folds when available and otherwise stratified
5-fold. Replica seeds are derived from a single \texttt{seed\_base}
(default 42) so the protocol can be replayed byte-for-byte.

\paragraph{Budget equivalence.} Single-objective + multi-objective
non-fidelity methods all run at the same headline budget
$B = \mathrm{DOE\_RUNS} + \mathrm{NBI\_EVAL\_K} = 88 + 50 = 138$
configurations per replica per (task, algorithm), each evaluated
under the same CV split. Multi-fidelity methods are budgeted by
total boosting iterations rather than configuration count
($\text{total\_n\_estimators\_budget} = B \cdot \texttt{max\_iter}$).
Evolutionary methods set population size and generations so
$\text{pop\_size} \cdot \text{gens} = B$. The full specification is in
Section~``Budget equivalence rule'' of
\texttt{docs/COMPARATIVE\_PROTOCOL.md}.

\paragraph{Headline target.} For the methods that run on all 72 CC18
tasks (every row of the method matrix with \texttt{full\_cc18}~$=$
\texttt{true}), the campaign is

\[
  72\ \text{tasks} \times 3\ \text{algorithms} \times
  n_{\text{methods,full}} \times 30\ \text{replicas},
\]
plus the per-replica jobs contributed by the \emph{subset} methods on
their defined CC18 subset. The exact $n_{\text{methods,full}}$ is
counted from \texttt{method\_matrix.csv} at shard-generation time and
is not duplicated as a constant in the manuscript.

\subsection{Metrics and statistical analysis}
\label{sec:exp-metrics}
\paragraph{Quality metrics.} ROC-AUC, PR-AUC, F1 (binary) /
F1-macro (multiclass), balanced accuracy, MCC. Multiclass tasks use
the one-vs-rest variants of the AUC metrics.

\paragraph{Calibration metrics.} Brier score (binary or multiclass)
and Expected Calibration Error (ECE)
\citep{naeini2015calibration}.

\paragraph{Cost metrics.} Mean fit-plus-predict time per fold;
total optimization time; total wall-clock time per replica.

\paragraph{Multi-objective frontier diagnostics.} For
multi-objective methods we additionally report: number of unique
non-dominated points; spread $\Delta$; weight-concentration index;
NBI sub-problem residual statistics (where applicable); MBPA-trigger
state. These are emitted regardless of whether MBPA fired so that
the conditional decision is always auditable.

\paragraph{Statistical analysis.} Paired comparisons within each task
use a non-parametric Wilcoxon signed-rank test (or paired $t$ after
log-transform when the residuals are reasonably Gaussian). Across
tasks, we use the Friedman test followed by Nemenyi post-hoc with
critical-difference plots, following
\citet{demsar2006statistical}. All $p$-values use Holm--Bonferroni
correction within each family of comparisons. Effect sizes are
reported as the median of paired differences plus a non-parametric
bootstrap 95\% confidence interval; rank statistics are computed
both per task and aggregated across the 72 tasks.

\paragraph{Robustness across replicas.} For each (task, algorithm,
method) cell we report the across-replica IQR of every quality
metric, and we additionally report the rate of replica-level
failures (job exits with status \texttt{failed}) so that the
robustness of each method is visible at the task level.

\subsection{Cost discipline and reproducibility}
\label{sec:exp-cost-discipline}
\paragraph{Cost discipline.} Before any shard generation,
\texttt{doe-xgb estimate-cost --preset
openml\_cc18\_72\_tasks\_3\_algorithms\_30\_replicas} produces an
authoritative wall-clock estimate. The campaign is gated by a
sign-off on this estimate.

\paragraph{Determinism.} Headline tables use
\texttt{tree\_method=exact} and \texttt{n\_jobs=1} (article-track
default); the determinism cost is documented in
\texttt{docs/METHODOLOGY\_DECISIONS.md} D13.

\paragraph{Sharded execution.} The job matrix is sharded across
SQLite files under \texttt{jobs/doctoral/openml\_cc18/shards/} so
that the dedicated MacBook Pro can run the campaign with low lock
contention; the schema (Section~\ref{sec:method-infra}) tracks
claim / success / failure transitions and per-job runtime so the
campaign can resume after interruption without re-running successful
jobs.

\paragraph{Reproducibility.} Every replica writes a
system-fingerprinted manifest, the full Pareto front (for
multi-objective methods), the NBI sub-problem residuals
(for the proposed method), the MBPA trigger diagnostics, and a
long-format per-fold metric CSV. Aggregation scripts regenerate
Tables~\ref{tab:performance}--\ref{tab:cost} of
Section~\ref{sec:results} from those artifacts.
